% block_mean.tex -- round 442, the unconditional block mean
% of kappa^dagger reduced to a frequency of q^dagger > 1
% (reviewer-R439 after r440).
% REDUZIERT / DICTIONARY_EXACT /
% UNSIGNED_MAJORANT_REFUTED / SELECTED_POINTWISE_CENSUS /
% MEAN_BOUND_OPEN.
% Builds with pdflatex.
% Companion machine check: verify_block_mean.py.
% Discovery: experiments/tfpt-discovery/block_mean_probe.py.
% NO RH CLAIM.  Research documentation, not a theorem of RH.
\documentclass[11pt]{article}

\usepackage[a4paper,margin=2.7cm]{geometry}
\usepackage{amsmath,amssymb,amsthm}
\usepackage{microtype}
\usepackage{booktabs}
\usepackage{xcolor}
\usepackage[colorlinks=true,linkcolor=blue!50!black,urlcolor=blue!50!black]{hyperref}

\newtheorem{lemma}{Lemma}
\newtheorem{prop}{Proposition}
\theoremstyle{remark}
\newtheorem{remark}{Remark}

\newcommand{\Rd}{R^{\dagger}}
\newcommand{\qd}{q^{\dagger}}
\newcommand{\kapd}{\kappa^{\dagger}}
\newcommand{\taud}{\tau^{\dagger}}
\newcommand{\Gd}{G^{\dagger}}
\newcommand{\Td}{\mathfrak{T}^{\dagger}}
\newcommand{\indm}{\operatorname{ind}_{-}}
\newcommand{\half}{\tfrac12}

\title{\bfseries The unconditional block mean\\[2mm]
\large one lemma: $\kapd=1\{\qd>1\}$;\\
the mean bound stays open}
\author{Stefan Hamann}
\date{August 30, 2026}

\begin{document}
\maketitle

\begin{abstract}\noindent
r440 showed that the block mean of
$\kapd=\indm(\Rd-\half I)$ is the winding of an averaged
integrand, and that a fixed contour about $(0,1)$ sees
none of $\kapd$ (the $1/2$-cluster sits in a collar of
$s=1$).  Combined with r433, that collar correction is
the last Redheffer pivot: $\kapd=1\{\qd>1\}$.
The block mean of $\kapd$ is therefore the frequency of
$\qd>1$ on the block.  This is the one lemma of the
round.  An unconditional bound keeping that frequency
strictly below $1$ is \emph{not} claimed: the unsigned
Chebyshev/Mertens majorant already exceeds $1$, and a
signed-oscillatory bound would sit at the circularity
gate.  No RH claim.
\end{abstract}

\medskip\noindent
\textbf{Status.}\enspace
Lemma~\ref{lem:dict} is \textbf{SATZ}.
Propositions~\ref{prop:src}--\ref{prop:unsigned} are
machine-checked identities or named refutations.
Proposition~\ref{prop:census} is a sealed census.
The unconditional frequency bound is \emph{open}.
Ausgang
\texttt{REDUZIERT / DICTIONARY\_EXACT / UNSIGNED\_MAJORANT\_REFUTED / SELECTED\_POINTWISE\_CENSUS / MEAN\_BOUND\_OPEN}.
No $L^{*}$ claim.  No $R^{\dagger}$ claim.  No RH claim.

\section{The lemma}\label{sec:lemma}

r440 T1: $\kapd=\#\{s\in(0,1):\taud(s)=0\}$ with
$\taud(s)=\det(I-s\Gd)$, $\Gd=(\Td)^{*}\Td$.
A circle of radius $0.40$ about $1/2$ winds $0$ on every
living window measured --- the interior $(0.10,0.90)$ is
empty.  So $\kapd$ is entirely the collar correction:
which side of $s=1$ the one collar zero occupies.
r433 identifies that side with the last Redheffer pivot
$\delta=1-\qd=-\mathrm{sch}$.

\begin{lemma}[Block-mean dictionary]\label{lem:dict}
On the r409/r433 construction,
$\kapd=1\{\qd>1\}$.
Hence on any finite block,
\[
\frac1K\sum_{k}\kapd_{k}
=\frac1K\sum_{k}1\{\qd_{k}>1\}.
\]
\end{lemma}

\begin{proof}
Spectral calculus plus the last-pivot identity.
Over $\mathbb{Q}$: $\delta=187/450$, $\qd=263/450<1$,
$1-\qd=-\mathrm{sch}$.
On w9 the three counters
($\indm(\Rd-\half I)$, zeros of $\taud$ in $(0,1)$,
$1\{\qd>1\}$) agree at $0$, with
$|1-\qd+\mathrm{sch}|=9.5\cdot 10^{-13}$.
On selected $k=3,4,5,6,7,9$ they agree at $0$;
on dead $\chi$ (six windows) they agree at $1$.
\end{proof}

r430 \texttt{exists\_index\_zero\_of\_block\_mean\_lt\_one}
then turns $\liminf$ of the right-hand side $<1$ into an
index-zero window.  The remaining analytic task is a bound
on the frequency, not a new identity.

\section{Source form of $\qd$}\label{sec:src}

r420/r424 write
$\qd=\gamma+v^{T}(I-BB^{T})^{-1}v$ with
$\gamma=\lVert b\rVert^{2}/B_{w}$ and $b$ the $\mu$-OP
coefficient of the signed border measure $\sigma$.
Woodbury gives the compact form
$\qd=B_{w}^{-1}\,b^{T}(I-C)^{-1}b$, $C=B^{T}B$.

\begin{prop}[Double sum]\label{prop:src}
\[
\qd=\sum_{a,a'}\frac{\sigma_{a}\sigma_{a'}}{B_{w}}
\bigl\langle\phi(t_{a}),(I-C)^{-1}\phi(t_{a'})\bigr\rangle.
\]
On w9 the two reconstructions match $\qd$ at
$1.2\cdot 10^{-15}$.
\end{prop}

Selected anchors $a_{k}=2^{k}$ are $\zeta$-windows: the
atoms $a$ include primes in dyadic ranges.  The
block-averaged $\qd$ is therefore an average of these
finite double sums.  Markov's inequality would convert
$\frac1K\sum\qd_{k}<1$ into a positive density of
$\qd<1$, hence into frequency $<1$.  That implication is
trivial; the missing input is the bound on the mean.

\section{Gate: unsigned majorants fail}\label{sec:gate}

\begin{prop}[Unsigned envelope]\label{prop:unsigned}
Replacing $b$ by $|b|$ on w9 yields $38.49\gg 1$.
Replacing $\sigma$ by $|\sigma|$ yields $2.10$ already
at $k=5$ ($k_{z}=17$) and $87.1$ at $k=9$ ($k_{z}=116$).
\end{prop}

A positive Chebyshev/Mertens majorant on the dyadic
prime mass therefore cannot prove $\qd<1$.
The remaining bound needs the \emph{signed} pairing of
$\sigma$ against $(I-C)^{-1}$.
PNT controls the main term of $\Lambda$ but not that
pairing.  An RH-quality error term on the oscillatory
sum is the circularity gate (forbidden as a mean-value
carrier).  The unconditional-positive route is closed.

\section{Census and trend}\label{sec:census}

\begin{prop}[Selected liveness]\label{prop:census}
Selected $k=3,4,5,6,7,9$: $\qd\in[0.891,0.962]$, all
$<1$, mean $0.932$, $\kapd=0$ on $6/6$
($k=8$ not rebuilt).
The sequence is \emph{not} falling: dip at $k=5$ then
rise toward $1$ ($k=9$: $0.962$).  Markov room $0.068$.
EXT-6: all $\qd<1$, mean $0.964$.
Core-42: $42/42$ $\qd<1$, mean $0.961$,
$\min=0.908$ ($k_{z}22$), $\max=0.9985$ ($k_{z}16$);
four windows sit above $0.99$ and stay living.
Dead $\chi$ $6/6$: $\qd>1$ and $\kapd=1$.
Living $\chi_{3}$-$9$: $\qd=0.935$, $\kapd=0$.
No dead Selected window is known.
\end{prop}

The selected sequence is a $\zeta$-window ladder.
r428 already recorded that the old ladder-criticality
was a quantifier artefact; the present census is
compatible with Selected being pointwise living on the
measured range.  The honest risk is cofinal approach
of $\qd$ to $1$ from below (vacuous MAIN $k_{z}16$ sits
at $0.9985$), not a known crossing.

\section{Kills}\label{sec:kills}

\begin{itemize}
\item \textbf{Soft mutants.}
$\operatorname{tr}\Gd$ is $50.9$ on living w9 and $53.0$
on dead $\chi_{3}$-$15$; $M_{2}=16\operatorname{tr}((I-\Rd)^{4})$
is $32.96$ vs $39.01$.  Both $\gg$ the threshold
(r398, r440).
\item \textbf{MI2.}
On the block $\{w9,k_{z}5\}$ the averaged-integrand
winding equals the mean of windings at $4\cdot 10^{-15}$,
and the fixed $r=0.40$ contour winds $0$ (the collar
carries all of $\kapd$).
\item \textbf{Scramble / circularity.}
Unaugmented $\mathrm{nneg}(R)=21$.  A pointwise bound
$\lvert\sum\Lambda n^{-1/2-it}\rvert\ll X^{\varepsilon}$
is not a mean-value carrier.
\end{itemize}

\section*{Claim boundary}

Nothing in this note is evidence for or against the
Riemann hypothesis, in either direction.
Lemma~\ref{lem:dict} is finite linear algebra.
The frequency bound that would feed r430 is open.

\begin{thebibliography}{9}
\bibitem{r398} Round 398, \texttt{high\_moment\_inertia}.
\bibitem{r420} Round 420, \texttt{cj\_sigma}:
  $\mathrm{den}=1+\gamma-v^{T}s$.
\bibitem{r424} Round 424, \texttt{gamma\_chain}:
  $b$ is the $\mu$-Parseval coefficient of $\sigma$.
\bibitem{r428} Round 428, \texttt{qn\_reopened}:
  ladder criticality as a quantifier artefact.
\bibitem{r430} \texttt{exists\_index\_zero\_of\_block\_mean\_lt\_one}.
\bibitem{r433} Round 433, \texttt{edge\_redheffer}:
  $\delta=1-\qd=-\mathrm{sch}$.
\bibitem{r440} Round 440, \texttt{mean\_tau\_index}:
  T1/T2/MI2 and the collar anatomy.
\end{thebibliography}

\end{document}
